\begin{tikzpicture}[
  font=\sffamily,
  srcbox/.style={
    rectangle, rounded corners=3pt, draw=ceindigo!80, line width=0.8pt,
    fill=white, text width=6.05cm, align=left, inner sep=6pt, font=\sffamily\footnotesize
  },
  hdr/.style={
    rectangle, rounded corners=3pt, fill=ceindigo, text=white,
    font=\sffamily\bfseries\small, text width=6.05cm, align=center, inner sep=6pt
  },
  hdrExt/.style={
    rectangle, rounded corners=3pt, fill=ceteal, text=white,
    font=\sffamily\bfseries\small, text width=6.05cm, align=center, inner sep=6pt
  }
]
  % Columna izquierda: 4 fuentes internas usando posicionamiento relativo
  \node[hdr] (hi) at (0, 0) {A. FUENTES INTERNAS (Sector / Empresa)};
  \node[srcbox, below=0.22cm of hi] (i1)
    {\textcolor{cecrimson}{\bfseries 1. Lo inesperado:} éxito, fracaso o suceso imprevisto. Menor riesgo y mayor accesibilidad.};
  \node[srcbox, below=0.20cm of i1] (i2)
    {\textcolor{cecrimson}{\bfseries 2. La incongruencia:} brecha entre lo que la industria asume y lo que el cliente vive.};
  \node[srcbox, below=0.20cm of i2] (i3)
    {\textcolor{cecrimson}{\bfseries 3. Necesidad de proceso:} cuello de botella operativo que exige un eslabón técnico nuevo.};
  \node[srcbox, below=0.20cm of i3] (i4)
    {\textcolor{cecrimson}{\bfseries 4. Cambios en industria o mercado:} rupturas rápidas de estructura que descolocan al líder.};

  % Columna derecha: 3 fuentes externas + síntesis
  \node[hdrExt, right=0.55cm of hi] (he) {B. FUENTES EXTERNAS (Entorno Social)};
  \node[srcbox, draw=ceteal!80, below=0.22cm of he] (e1)
    {\textcolor{ceteal}{\bfseries 5. Cambios demográficos:} edad, renta, urbanización o migraciones. Altamente predecibles.};
  \node[srcbox, draw=ceteal!80, below=0.20cm of e1] (e2)
    {\textcolor{ceteal}{\bfseries 6. Percepción y significado:} giro cultural en hábitos de salud, sostenibilidad o tiempo.};
  \node[srcbox, draw=ceteal!80, below=0.20cm of e2] (e3)
    {\textcolor{ceteal}{\bfseries 7. Nuevo conocimiento:} ciencia, IA o patentes. Alto impacto, mayor plazo y mayor riesgo.};
  \node[rectangle, rounded corners=3pt, fill=cepanelamber, draw=ceamber, line width=0.9pt,
        text width=6.05cm, align=left, inner sep=6pt, font=\sffamily\footnotesize,
        below=0.20cm of e3] (nota) {
    \textcolor{ceamber!90!black}{\bfseries Clave de Drucker (1985):} la innovación sistemática monitoriza estas 7 fuentes con método.
  };
\end{tikzpicture}